\subsection{Reactions and Rate Equations}
%TODO: Add references to the claims
The SCR-ASC system involves urea decomposition, ammonia adsorption/desorption, selective catalytic reduction of $\nox$,
and ammonia oxidation (AMOX) on both SCR and ASC surfaces. The following assumptions reduce the full reaction network to
a tractable form.

\itbf{$A_1$:} All $\nox$ in the exhaust is $\nitrox$, since commercially available $\nox$ sensors cannot differentiate
between $\nitrox$ and $\notwo$ \cite{nova2014urea}.

\itbf{$A_2$:} Three dominant reactions are retained: the standard SCR reaction, ammonia oxidation (AMOX), and ammonia
adsorption/desorption. Fast and slow SCR reactions are omitted, as the exhaust flow rate ensures they do not
significantly affect tailpipe concentrations.

\itbf{$A_3$:} The AMOX reactions on both SCR and ASC surfaces are lumped into a single oxidation reaction with $100\%$
$\nitrogen$ selectivity (valid for $T > 225\,°\text{C}$).

\itbf{$A_4$:} The Eley-Rideal mechanism governs the standard SCR reaction: gaseous $\nox$ reacts with surface-adsorbed
$\ammonia$.

\itbf{$A_5$:} Mass transfer is neglected; the catalyst kinetics are reaction-controlled, as the standard SCR reaction
rate exceeds the exhaust fluid flow rate.

\itbf{$A_6$:} The SCR-ASC chamber is modelled as a plug-flow reactor; reaction rates within each residence time are
governed by inlet gas-phase concentrations.
%===
\begin{enumerate}
    \item Standard SCR reaction:
          \begin{align}
              4 \ammonia^{ads} + 4 \nitrox + \oxygen & \xrightarrow[]{k_{scr}} 4 \nitrogen + 6 \water \label{react::std_scr}
          \end{align}
    \item Ammonia Oxidation:
          \begin{align}
              4 \ammonia^{ads} + 3 \oxygen & \xrightarrow[]{k_{oxi}} 2 \nitrogen + 6 \water \label{react::amox}
          \end{align}
    \item Ammonia Adsorption/Desorption:
          \begin{align}
              \ammonia + \Theta_{free} & \xrightleftharpoons[k_{des}]{k_{ads}} \ammonia^{ads}
              \label{react::ads}
          \end{align}
\end{enumerate}
%===
The surface $\ammonia$ rate follows directly from these three reactions. Adsorption fills sites proportional to
free-site concentration $(\Gamma - \con{\ammonia}^{ads})$, while desorption, SCR, and oxidation each consume adsorbed
$\ammonia$:
\begin{multline}
    \dot{\con{\ammonia}}^{ads} = \underbrace{k_{ads}\,\con{\ammonia}^{in}\,\lr{\Gamma - \con{\ammonia}^{ads}}}_{\text{adsorption}} \\
    - \underbrace{\lr{k_{des} + k_{scr}\,\con{\nox}^{in} + k_{oxi}}\,\con{\ammonia}^{ads}}_{\text{desorption + SCR + oxidation}}
    \label{eqn::nh3_ads_rate}
\end{multline}
where $\Gamma$ and $\con{\ammonia}^{ads}$ are surface concentrations ($mol/m^2$), while the remaining quantities are
volumetric concentrations ($mol/m^3$). The conversion between surface and volumetric domains is given by the
surface-to-volume ratio $k_{s2v} = A_{scr}/V$ ($m^{-1}$). The volumetric $\nox$ consumption rate involves only the SCR
reaction:
\begin{align}
    \dot{\con{\nox}}^{scr} &= -k_{s2v}\,k_{scr}\,\con{\ammonia}^{ads}\,\con{\nox}^{in}
    \label{eqn::nox_rate}
\end{align}
These two rate equations are coupled through $\con{\ammonia}^{ads}$. The surface concentration drives $\nox$
reduction but is itself unobservable from tailpipe measurements. Eliminating this hidden state and expressing the $\nox$
dynamics purely in terms of measurable signals is the central objective of the derivation that follows.